\section{A Laplace lemma with a parameter, and the next term}\label{app:laplace}

This appendix proves the Laplace lemma that Appendix~\ref{app:torus} uses
for Theorem~\ref{thm:local}, and then Proposition~\ref{prop:thirdorder}. Let
$\mathbb T^m=(\mathbb R/2\pi\mathbb Z)^m$. For $\varphi\in\mathbb T^m$ we use
the max-norm $\lVert\varphi\rVert=\max_i\operatorname{dist}(\varphi_i,2\pi\mathbb Z)$.
For $x\in\mathbb R^m$ we write $\lVert x\rVert=\max_i|x_i|$ and $|x|$ for the
Euclidean norm, so that $\lVert x\rVert\le|x|$. Repeated indices are summed
from $1$ to $m$.

\begin{lemma}[Laplace integrals with a parameter]\label{lem:laplace}
Let $\mathcal K$ be a compact metric space and $m\ge1$. Let
$\Psi,a:\mathcal K\times\mathbb T^m\to\mathbb C$ be such that the
$\varphi$-derivatives of $\Psi$ up to order $5$ and of $a$ up to order $3$
exist and are jointly continuous on $\mathcal K\times\mathbb T^m$. Assume that
for every $\alpha\in\mathcal K$
\begin{enumerate}[(i)]
\item $\Psi(\alpha,0)\in\mathbb R$ and $\nabla_\varphi\Psi(\alpha,0)=0$,
\item $\mathcal J(\alpha)=-\nabla^2_\varphi\Psi(\alpha,0)$ is real, symmetric
and positive definite, and
\item $\operatorname{Re}\Psi(\alpha,\varphi)<\Psi(\alpha,0)$ for every
$\varphi\ne0$ in $\mathbb T^m$.
\end{enumerate}
Then, uniformly for $\alpha\in\mathcal K$ and $T\ge1$,
\[
\int_{\mathbb T^m}e^{T\Psi(\alpha,\varphi)}a(\alpha,\varphi)\,d\varphi
=e^{T\Psi(\alpha,0)}\Bigl(\frac{2\pi}T\Bigr)^{m/2}\det\mathcal J(\alpha)^{-1/2}
\Bigl(a_0+\frac{c_1}T+O(T^{-3/2})\Bigr).
\]
Here $a_0=a(\alpha,0)$. Write $\mathcal J^{ij}$ for the entries of
$\mathcal J(\alpha)^{-1}$, and $\Psi_{ijk}$, $\Psi_{ijkl}$, $a_i$, $a_{ij}$ for
the $\varphi$-derivatives at $(\alpha,0)$. Then
\[
c_1=\tfrac12a_{ij}\mathcal J^{ij}
+\tfrac12a_i\Psi_{jkl}\mathcal J^{ij}\mathcal J^{kl}
+a_0\Bigl[\tfrac18\Psi_{ijkl}\mathcal J^{ij}\mathcal J^{kl}
+\tfrac1{12}\Psi_{ijk}\Psi_{lrs}\mathcal J^{il}\mathcal J^{jr}\mathcal J^{ks}
+\tfrac18\Psi_{ijk}\Psi_{lrs}\mathcal J^{ij}\mathcal J^{kl}\mathcal J^{rs}\Bigr].
\]
It is continuous on $\mathcal K$, and it is real-analytic in $\alpha$ when
$\Psi$ and $a$ are real-analytic in $(\alpha,\varphi)$.
\end{lemma}

\begin{proof}
\emph{Constants.} The derivatives are jointly continuous on the compact set
$\mathcal K\times\mathbb T^m$. So
$B_\Psi=\max_{|\beta|\le5}\sup|\partial^\beta_\varphi\Psi|$ and
$B_a=\max_{|\beta|\le3}\sup|\partial^\beta_\varphi a|$ are finite, and
$\mathcal J(\alpha)$ is continuous in $\alpha$. Hence the least eigenvalue of
$\mathcal J(\alpha)$ has a positive minimum $h_{\min}$ over $\mathcal K$.
Also $\det\mathcal J(\alpha)\le(mB_\Psi)^m$, since every eigenvalue is at
most the largest absolute row sum. Below, $C$ is a constant that depends only
on $m$, $h_{\min}$, $B_\Psi$, $B_a$ and the numbers $\varepsilon$ and $\eta$
chosen below. It may change from line to line.

\emph{The chart.} Let $C_3=m^3B_\Psi/6$ and choose $\varepsilon$ with
$0<\varepsilon<1$ and $C_3\varepsilon\le h_{\min}/4$. Since $\varepsilon<\pi$,
reduction modulo $2\pi$ is a measure-preserving bijection from the cube
$\{x\in\mathbb R^m:\lVert x\rVert\le\varepsilon\}$ onto
$\{\varphi\in\mathbb T^m:\lVert\varphi\rVert\le\varepsilon\}$, and we identify
the two. The cube is convex, so Taylor's theorem with integral remainder
applies to $\Psi(\alpha,\cdot)$ and $a(\alpha,\cdot)$ on it. It holds for
complex values, since the remainder is an integral. Put
$R_3(\alpha,\varphi)=\Psi(\alpha,\varphi)-\Psi(\alpha,0)
+\frac12\varphi^\top\mathcal J(\alpha)\varphi$. By (i) and (ii), $R_3$ and its
derivatives of order $1$ and $2$ vanish at $\varphi=0$. The remainder formula
of order $3$ and $\sum_{|\beta|=3}1/\beta!=m^3/6$ give
$|R_3(\alpha,\varphi)|\le C_3\lVert\varphi\rVert^3$. Since
$\lVert\varphi\rVert^3\le\lVert\varphi\rVert\,|\varphi|^2$, we get
\begin{equation}\label{eq:laplace-R3}
|R_3(\alpha,\varphi)|\le C_3\varepsilon|\varphi|^2\le\tfrac14h_{\min}|\varphi|^2
\qquad\text{for }\lVert\varphi\rVert\le\varepsilon .
\end{equation}

\emph{Step 1 (localization).} The set of $(\alpha,\varphi)$ with
$\lVert\varphi\rVert\ge\varepsilon$ is compact, and
$\operatorname{Re}\Psi(\alpha,\varphi)-\Psi(\alpha,0)$ is continuous and, by
(iii), negative on it. So it is at most $-\eta$ there, for some $\eta>0$. Hence
the integral over $\lVert\varphi\rVert\ge\varepsilon$ has modulus at most
$(2\pi)^mB_ae^{T(\Psi(\alpha,0)-\eta)}$. Since
$T^{m/2+3/2}e^{-\eta T}\le C$, this is at most
$CT^{-3/2}\cdot e^{T\Psi(\alpha,0)}T^{-m/2}$.

\emph{Step 2 (scaling).} On $\lVert\varphi\rVert\le\varepsilon$ we substitute
$\varphi=x/\sqrt T$. With $w(x)=TR_3(\alpha,x/\sqrt T)$, the integral over
this set is
\[
T^{-m/2}e^{T\Psi(\alpha,0)}\int_{\lVert x\rVert\le\varepsilon\sqrt T}
e^{-\frac12x^\top\mathcal Jx+w(x)}\,a(\alpha,x/\sqrt T)\,dx .
\]
By~\eqref{eq:laplace-R3}, $|w(x)|\le\frac14h_{\min}|x|^2$. Since
$x^\top\mathcal Jx\ge h_{\min}|x|^2$, we get
$|e^{-\frac12x^\top\mathcal Jx+w(x)}|\le e^{-h_{\min}|x|^2/4}$.

\emph{Step 3 (expansion).} Put
$\mathcal P_3(x)=\frac16\Psi_{ijk}x^ix^jx^k$ and
$\mathcal P_4(x)=\frac1{24}\Psi_{ijkl}x^ix^jx^kx^l$. For
$\lVert x\rVert\le\varepsilon\sqrt T$, Taylor's theorem to order $5$ for
$\Psi$ and to order $3$ for $a$ gives
\[
w=T^{-1/2}\mathcal P_3+T^{-1}\mathcal P_4+\omega_1,\qquad
a(\alpha,x/\sqrt T)=a_0+T^{-1/2}a_ix^i+\tfrac1{2T}a_{ij}x^ix^j+\omega_2,
\]
with $|\omega_1|\le CT^{-3/2}|x|^5$ and $|\omega_2|\le CT^{-3/2}|x|^3$. Also
$|w|\le C_3T^{-1/2}\lVert x\rVert^3$ and $e^{|w|}\le e^{h_{\min}|x|^2/4}$. Let
$\mathcal Y=1+T^{-1/2}\mathcal P_3+T^{-1}(\mathcal P_4+\frac12\mathcal P_3^2)$.
Then $e^w-\mathcal Y$ is the sum of $\omega_1$, of
$\frac12(w-T^{-1/2}\mathcal P_3)(w+T^{-1/2}\mathcal P_3)$ with
$w-T^{-1/2}\mathcal P_3=T^{-1}\mathcal P_4+\omega_1$, and of
$e^w-1-w-\frac12w^2$, which is at most $\frac16|w|^3e^{|w|}$ in modulus. So
$|e^w-\mathcal Y|\le CT^{-3/2}(1+|x|)^9e^{h_{\min}|x|^2/4}$. We multiply
$e^w=\mathcal Y+(e^w-\mathcal Y)$ by the expansion of $a$, and use $T\ge1$
for the products that carry $T^{-3/2}$ or $T^{-2}$. This gives
\[
e^{w}a(\alpha,x/\sqrt T)=a_0+T^{-1/2}\mathcal Q_1(x)+T^{-1}\mathcal Q_2(x)
+\omega(x),\qquad
|\omega(x)|\le CT^{-3/2}(1+|x|)^{11}e^{h_{\min}|x|^2/4},
\]
where $\mathcal Q_1=a_0\mathcal P_3+a_ix^i$ is odd and
$\mathcal Q_2=a_0(\mathcal P_4+\frac12\mathcal P_3^2)+a_ix^i\mathcal P_3
+\frac12a_{ij}x^ix^j$.

\emph{Step 4 (integration).} Since
$\frac12x^\top\mathcal Jx\ge\frac12h_{\min}|x|^2$, the remainder $\omega$
contributes at most
$CT^{-3/2}\int_{\mathbb R^m}(1+|x|)^{11}e^{-h_{\min}|x|^2/4}dx=CT^{-3/2}$ to the
integral in Step 2. If $\lVert x\rVert>\varepsilon\sqrt T$ then
$|x|^2>\varepsilon^2T$. So extending the integrals of the 3 polynomial terms
from $\lVert x\rVert\le\varepsilon\sqrt T$ to $\mathbb R^m$ changes them by at
most $Ce^{-h_{\min}\varepsilon^2T/4}\le CT^{-3/2}$. On $\mathbb R^m$,
$\int e^{-\frac12x^\top\mathcal Jx}dx=(2\pi)^{m/2}\det\mathcal J^{-1/2}$, and the
odd term $\mathcal Q_1$ integrates to $0$. This is why the first correction is
of order $1/T$ and not $T^{-1/2}$. For $\mathcal Q_2$ let $x$ be a centered
Gaussian vector with $\mathbb E[x^ix^j]=\mathcal J^{ij}$. By Isserlis' theorem,
a moment of order $4$ or $6$ is the sum, over the $3$ or $15$ ways to pair the
indices, of the products of the $\mathcal J^{ij}$ of the pairs. Since the
derivatives of $\Psi$ are symmetric in their indices, this gives
$\mathbb E[a_ix^i\mathcal P_3]=\frac12a_i\Psi_{jkl}\mathcal J^{ij}\mathcal J^{kl}$
and $\mathbb E[\mathcal P_4]=\frac18\Psi_{ijkl}\mathcal J^{ij}\mathcal J^{kl}$.
In
$\mathbb E[\mathcal P_3^2]=\frac1{36}\Psi_{ijk}\Psi_{lrs}
\mathbb E[x^ix^jx^kx^lx^rx^s]$, the $6$ pairings that join each index of the
first factor to an index of the second each give
$\Psi_{ijk}\Psi_{lrs}\mathcal J^{il}\mathcal J^{jr}\mathcal J^{ks}$. The other
$9$ pair 2 indices inside each factor and join the remaining 2, and each gives
$\Psi_{ijk}\Psi_{lrs}\mathcal J^{ij}\mathcal J^{kl}\mathcal J^{rs}$. So
$\mathbb E[\mathcal Q_2]=c_1$.

By Steps 2 to 4 the integral over $\lVert\varphi\rVert\le\varepsilon$ is
$e^{T\Psi(\alpha,0)}T^{-m/2}\bigl[(2\pi)^{m/2}\det\mathcal J^{-1/2}(a_0+c_1/T)
+O(T^{-3/2})\bigr]$. Since $\det\mathcal J\le(mB_\Psi)^m$, this error and the
one in Step 1 are $O(T^{-3/2})$ relative to
$e^{T\Psi(\alpha,0)}(2\pi/T)^{m/2}\det\mathcal J^{-1/2}$. This proves the
expansion. All constants depend only on $m$, $h_{\min}$, $\varepsilon$,
$\eta$, $B_\Psi$ and $B_a$, so it holds uniformly on $\mathcal K$. Finally
$c_1$ is a polynomial in $a_0$, the derivatives at $\varphi=0$ and the entries
of $\mathcal J^{-1}$, which gives its continuity and analyticity.
\end{proof}

\begin{proof}[Proof of Proposition~\ref{prop:thirdorder}]
If $k=1$, then $f_T=\mu_ie^{-q_iT}=K_Fe^{-T\lambda_F}$ exactly
(Theorem~\ref{thm:local}), so the bracket in (i) vanishes and $c_F=0$. This
proves (i), and (iii) for single states. Until the proof of (iii) let
$k\ge2$. We use the coordinates $\alpha'$ and the Euclidean norm
$|\cdot|$ of Section~\ref{sec:secondorder}, and put
$x=\alpha'-\alpha^{*\prime}_F$, $c^*=c(\alpha^*_F)$ and
$\bar c=c^*+\frac12\gamma^\top\Sigma\gamma$.

\emph{Taylor bounds.} By Proposition~\ref{prop:constant}(a) and
Theorem~\ref{thm:local}, the functions $I_F$, $\log C_F$ and $c$ are
real-analytic on $\Delta_F^\circ$. Here $C_F>0$ because $\mu(F)>0$. Also
$\nabla I_F(\alpha^*_F)=0$, since $\alpha^*_F$ is an interior minimum of $I_F$
(Proposition~\ref{prop:rate}(a)), and $\nabla^2I_F(\alpha^*_F)=\Sigma^{-1}$.
Let $4c_U>0$ be the least eigenvalue of $\Sigma^{-1}$. Then there are a closed
ball $U=\{|x|\le\rho_0\}\subset\Delta_F^\circ$ and $L_1>0$ with
$L_1\rho_0\le c_U$ such that for all $\alpha,\beta\in U$
\begin{gather*}
\Bigl|\log\frac{C_F(\alpha)}{K_F}-\gamma\cdot x\Bigr|\le L_1|x|^2,\qquad
\Bigl|I_F(\alpha)-\lambda_F-\tfrac12x^\top\Sigma^{-1}x\Bigr|\le L_1|x|^3,\\
|\log C_F(\alpha)-\log C_F(\beta)|+|c(\alpha)-c(\beta)|
+|\nabla I_F(\alpha)-\nabla I_F(\beta)|\le L_1|\alpha'-\beta'|,\\
|I_F(\alpha)-I_F(\beta)-\nabla I_F(\beta)\cdot(\alpha'-\beta')|
\le L_1|\alpha'-\beta'|^2 .
\end{gather*}
In particular $I_F(\alpha)-\lambda_F\ge2c_U|x|^2-L_1|x|^3\ge c_U|x|^2$ on $U$.

\emph{(i), outside $U$.} We first show that for $\theta\in V_F$, $T\ge1$ and
$\alpha\in\Delta_F^\circ$,
\begin{equation}\label{eq:contupper}
f_T(\alpha)\le C_\theta T^{k-1}e^{-T\langle q-\theta,\alpha\rangle},
\end{equation}
with $C_\theta$ from Lemma~\ref{lem:upper}. Fix $T$ and $\alpha$, and for large
$N$ choose $n$ with support $F$ and $|n_i-N\alpha_i|<1$ for all $i$. By
Lemma~\ref{lem:discrete} on a closed ball around $\alpha$ in
$\Delta_F^\circ$, and the continuity of $f_T$, we have
$N^{k-1}P_N(n)\to f_T(\alpha)$ as $N\to\infty$. By Lemma~\ref{lem:upper} with
$h=T/N$,
$N^{k-1}P_N(n)\le C_\theta T^{k-1}\exp(-T\langle q-\theta,n/N\rangle
+C_\theta T^2/N)$, and the right side tends to that of~\eqref{eq:contupper}.
The argument for the region outside $U$ in the proof of
Theorem~\ref{thm:facemax} gives $\eta>0$ and
$\theta_1,\dots,\theta_J\in V_F$ with
$\max_j\langle q-\theta_j,\alpha\rangle>\lambda_F+2\eta$ on
$\Delta_F\setminus U^\circ$. So~\eqref{eq:contupper} gives
$f_T(\alpha)\le CT^{k-1}e^{-T(\lambda_F+2\eta)}$ for
$\alpha\in\Delta_F^\circ\setminus U$.

\emph{(i), inside $U$.} Put $X_T=T^{(k-1)/2}K_Fe^{-T\lambda_F}$ and
$R_T=f_T/X_T$. By Theorem~\ref{thm:local} with $\mathcal K=U$, for
$\alpha\in U$,
\begin{equation}\label{eq:RT}
R_T(\alpha)=\frac{C_F(\alpha)}{K_F}e^{-T(I_F(\alpha)-\lambda_F)}
\Bigl(1+\frac{c(\alpha)}T+\epsilon_T(\alpha)\Bigr),\qquad
|\epsilon_T(\alpha)|\le C_1T^{-3/2}.
\end{equation}
Let $x_0=\Sigma\gamma/T$, and for every $x$ put $y=x-x_0$. Completing the
square gives
\begin{equation}\label{eq:square}
\gamma\cdot x-\frac T2x^\top\Sigma^{-1}x
=\frac{\gamma^\top\Sigma\gamma}{2T}-\frac T2y^\top\Sigma^{-1}y .
\end{equation}
For large $T$ the point $\alpha_0$ with $x=x_0$ lies in $U$. There the Taylor
bounds give
$\log(C_F/K_F)-T(I_F-\lambda_F)\ge\frac{\gamma^\top\Sigma\gamma}{2T}
-L_1|x_0|^2-L_1T|x_0|^3$ and $c(\alpha_0)\ge c^*-L_1|x_0|$. Since
$|x_0|=O(1/T)$,
\begin{equation}\label{eq:Rlower}
R_T(\alpha_0)\ge1+\frac{\bar c}T-O(T^{-3/2}).
\end{equation}
Now let $\alpha\in U$. If $|x|\ge T^{-1/3}$, then
$\log(C_F/K_F)-T(I_F-\lambda_F)\le|\gamma||x|+L_1|x|^2-c_UT|x|^2
\le-\frac{c_U}2T^{1/3}$ for large $T$, so $R_T(\alpha)\le2e^{-c_UT^{1/3}/2}$.
If $|x|\le T^{-1/3}$, then $|y|\le2T^{-1/3}$ for large $T$. From
$|x|^2\le2|y|^2+2|x_0|^2$ and $|x|^3\le4|y|^3+4|x_0|^3$ we get
$L_1|x|^2+L_1T|x|^3\le(2L_1+8L_1T^{2/3})|y|^2+O(T^{-2})\le c_UT|y|^2+O(T^{-2})$
for large $T$. With~\eqref{eq:square} and
$\frac T2y^\top\Sigma^{-1}y\ge2c_UT|y|^2$, this gives
\[
\log\frac{C_F(\alpha)}{K_F}-T(I_F(\alpha)-\lambda_F)
\le\frac{\gamma^\top\Sigma\gamma}{2T}-c_UT|y|^2+O(T^{-2}).
\]
Also $c(\alpha)\le c^*+L_1|y|+L_1|x_0|$ and
$\sup_{u\ge0}ue^{-c_UTu^2}=O(T^{-1/2})$. Since $1+c^*/T>0$ for large $T$,
\eqref{eq:RT} gives
\begin{equation}\label{eq:Rupper}
R_T(\alpha)\le e^{\gamma^\top\Sigma\gamma/(2T)+O(T^{-2})}
\Bigl[e^{-c_UT|y|^2}\Bigl(1+\frac{c^*}T\Bigr)+O(T^{-3/2})\Bigr]
\le1+\frac{\bar c}T+O(T^{-3/2}).
\end{equation}
Off $U$ we showed $R_T=O(T^{(k-1)/2}e^{-2\eta T})$. Together
with~\eqref{eq:Rlower} and~\eqref{eq:Rupper}, this shows that as
$T\to\infty$
\begin{equation}\label{eq:supf}
\sup_{\alpha\in\Delta_F^\circ}f_T(\alpha)=T^{(k-1)/2}K_Fe^{-T\lambda_F}
\Bigl(1+\frac{\bar c}T+O(T^{-3/2})\Bigr).
\end{equation}
So the limit in (i) exists and equals $\bar c$.

We also locate the maximizer. For large $T$, the value $R_T(\alpha_0)$ exceeds
the bounds above off $U$ and for $|x|\ge T^{-1/3}$. Since $f_T$ is continuous
on $U$, it has a maximizer $\hat\alpha$ over $\Delta_F^\circ$, with
$\hat\alpha\in U$ and $|\hat x|\le T^{-1/3}$. Comparing the middle
of~\eqref{eq:Rupper} at $\hat\alpha$ with~\eqref{eq:Rlower} gives
$e^{-c_UT|\hat y|^2}\ge1-O(T^{-3/2})$. So $|\hat y|=O(T^{-5/4})$ and
\begin{equation}\label{eq:maxloc}
|\hat\alpha'-\alpha^{*\prime}_F|\le|x_0|+|\hat y|=O(1/T).
\end{equation}

\emph{(ii).} We show that for every $F$ that satisfies (H) there are $T_0$
and $N_0$ such that, for $N\ge N_0$ and $T_0\le T\le N^{1/3}$,
\begin{equation}\label{eq:logmaxnext}
\log M_F=-(k-1)L+\frac{k-1}2\log T+\log K_F-T\lambda_F+\frac{c_F}T
+O\Bigl(T^{-3/2}+\frac{T^2}N\Bigr).
\end{equation}
For $k=1$ this is Theorem~\ref{thm:facemax}, since $c_F=0$. Let $k\ge2$ and
$X=N^{-(k-1)}X_T$. For $n/N\in U$, Lemma~\ref{lem:discrete}
and~\eqref{eq:supf} give
$P_N(n)\le N^{-(k-1)}\sup f_T\,(1+O(T^2/N))
=X(1+c_F/T+O(T^{-3/2}+T^2/N))$. For $n/N\notin U$, Lemma~\ref{lem:upper}
with the $\theta_j$ above and $T^2/N\le1$ gives
$P_N(n)\le C(T/N)^{k-1}e^{-T(\lambda_F+2\eta)}=X\cdot
O(T^{(k-1)/2}e^{-2\eta T})$.

For the lower bound let $\hat\alpha$ be the maximizer
in~\eqref{eq:maxloc}. Choose $n$ with support $F$ and
$|n_i-N\hat\alpha_i|<1$ for all $i$, and put $\beta=n/N$, so that
$|\beta'-\hat\alpha'|\le k/N$. For $T\ge T_0$ and $N\ge N_0$ both $\hat\alpha$
and $\beta$ lie in $U$, and we apply Theorem~\ref{thm:local} at both points.
By the Taylor bounds, $|\log C_F(\beta)-\log C_F(\hat\alpha)|$ and
$|c(\beta)-c(\hat\alpha)|$ are at most $L_1k/N$. Also
$|\nabla I_F(\hat\alpha)|=|\nabla I_F(\hat\alpha)-\nabla I_F(\alpha^*_F)|
\le L_1|\hat x|=O(1/T)$ by~\eqref{eq:maxloc}, so
\[
T|I_F(\beta)-I_F(\hat\alpha)|\le T\bigl(|\nabla I_F(\hat\alpha)|\,
|\beta'-\hat\alpha'|+L_1|\beta'-\hat\alpha'|^2\bigr)=O(1/N),
\]
since $T\le N$. Hence $f_T(\beta)=f_T(\hat\alpha)(1+O(1/N+T^{-3/2}))$. Note
that this step uses Theorem~\ref{thm:local} at the 2 points, and no bound on
the derivatives of $f_T$. By Lemma~\ref{lem:discrete} and~\eqref{eq:supf},
\[
M_F\ge P_N(n)=N^{-(k-1)}f_T(\beta)\bigl(1+O(T^2/N)\bigr)
=X\bigl(1+c_F/T-O(T^{-3/2}+T^2/N)\bigr),
\]
since $1/N\le T^2/N$. So $M_F=X(1+c_F/T+O(T^{-3/2}+T^2/N))$. Taking
logarithms gives~\eqref{eq:logmaxnext}, because $T^2/N\le N^{-1/3}$.

Now let $T_N\in[1,N^{1/3}]$ be a zero of $\log M_F-\log M_G$.
By~\eqref{eq:crossing}, $T_N\to\infty$, so $T_N\ge T_0$ for large $N$.
Subtracting~\eqref{eq:logmaxnext} for $G$ from~\eqref{eq:logmaxnext} for $F$
at $T=T_N$ gives
\[
\Delta L-\frac\Delta2\log T_N-\delta T_N-\log\frac{K_F}{K_G}
=\frac{c_F-c_G}{T_N}+O\Bigl(T_N^{-3/2}+\frac{T_N^2}N\Bigr),
\]
and multiplying by $T_N$ gives (ii).

\emph{(iii).} Single states were treated at the start. In the 2 other cases a
group $S$ of permutations of $F$ acts transitively on $F$ and preserves $A_F$,
$q$ and the restriction of $\mu$ to $F$. We first show that then $\alpha^*_F$
is uniform on $F$ and $\gamma=0$, so that $c_F=c(\alpha^*_F)$. For
$\sigma\in S$ and $v\in\mathbb R^F$ let $\sigma v$ be $v$ with its coordinates
permuted by $\sigma$. A permutation in $S$ maps each word in the definition of
$W$ to a word with the same weight, so $H(\sigma t)=H(t)$ and
$f_T(\sigma\alpha)=f_T(\alpha)$. Substituting $u\mapsto\sigma u$ in the
definition of $I_F$ gives $I_F(\sigma\alpha)=I_F(\alpha)$. So
$\sigma\alpha^*_F$ also minimizes $I_F$, and $\sigma\alpha^*_F=\alpha^*_F$ by
Proposition~\ref{prop:rate}(a). Since $S$ is transitive, $\alpha^*_F$ is
uniform on $F$. By Theorem~\ref{thm:local},
$C_F(\alpha)=\lim_{T\to\infty}T^{-(k-1)/2}e^{TI_F(\alpha)}f_T(\alpha)$, so
$C_F(\sigma\alpha)=C_F(\alpha)$ as well. Hence the derivative of $\log C_F$ at
$\alpha^*_F$ is a linear form on $\{v\in\mathbb R^F:\sum_iv_i=0\}$ that takes
the same value at $v$ and at $\sigma v$. The average of $\sigma v$ over
$\sigma\in S$ is the constant vector with entries $\frac1k\sum_iv_i=0$. So the
form vanishes and $\gamma=0$.

\emph{A symmetric pair.} Let $F=\{i,j\}$ with $q_{ij}=q_{ji}=\kappa$,
$q_i=q_j=q$ and $\mu_i=\mu_j=\mu_0$. By (H), $\kappa>0$ and $\mu_0>0$. The
swap of $i$ and $j$ gives $S$, so $\alpha^*_F=(\frac12,\frac12)$ and
$c_F=c(\alpha^*_F)$. Also $Q_F(1,1)^\top=(\kappa-q)(1,1)^\top$, so
$\lambda_F=q-\kappa$. The words in $F$ alternate. A word that starts at $i$
and contains $j$ has either $2l$ letters and run counts $(l,l)$, or $2l+1$
letters and run counts $(l+1,l)$, with $l\ge1$, and its weight is
$\mu_0\kappa^{M-1}$. The words that start at $j$ give the same terms at
$t_i=t_j$. With
$I_\nu(\xi)=\sum_{l\ge0}(\xi/2)^{2l+\nu}/(l!\,(l+\nu)!)$~\cite[Eq.~10.25.2]{dlmf}
we get
\[
H\Bigl(\frac T2,\frac T2\Bigr)=2\mu_0\sum_{l\ge1}
\Bigl(\frac{\kappa^{2l-1}(T/2)^{2l-2}}{((l-1)!)^2}
+\frac{\kappa^{2l}(T/2)^{2l-1}}{l!\,(l-1)!}\Bigr)
=2\mu_0\kappa\bigl(I_0(\kappa T)+I_1(\kappa T)\bigr).
\]
By Lemma~\ref{lem:continuum}, $f_T(\alpha^*_F)=Te^{-qT}H(T/2,T/2)$. As
$\xi\to\infty$,
$I_\nu(\xi)=(2\pi\xi)^{-1/2}e^\xi\bigl(1-\frac{4\nu^2-1}{8\xi}+O(\xi^{-2})\bigr)$
\cite[Eq.~10.40.1]{dlmf}, so
$I_0(\xi)+I_1(\xi)=2(2\pi\xi)^{-1/2}e^\xi\bigl(1-\frac1{8\xi}+O(\xi^{-2})\bigr)$.
Hence
\[
f_T(\alpha^*_F)=\frac{4\mu_0\sqrt\kappa}{\sqrt{2\pi}}\,T^{1/2}e^{-(q-\kappa)T}
\Bigl(1-\frac1{8\kappa T}+O(T^{-2})\Bigr).
\]
Theorem~\ref{thm:local} at $\alpha^*_F$ gives
$f_T(\alpha^*_F)=T^{1/2}K_Fe^{-(q-\kappa)T}(1+c(\alpha^*_F)/T+O(T^{-3/2}))$.
Letting $T\to\infty$ gives first $K_F=4\mu_0\sqrt\kappa/\sqrt{2\pi}$ and then
$c(\alpha^*_F)=-1/(8\kappa)$. So $c_F=-1/(8\kappa)$, which is $-\frac18$ when
$\kappa=1$. For $\kappa=1$ and $\mu_0=\frac12$ we recover
$K_F=\sqrt{2/\pi}$ from Corollary~\ref{cor:paperA}.

\emph{A face of the complete graph.} Let $Q=J-dI$, so that $q_{ij}=1$ for
$i\ne j$ and $q_i=d-1$. Let $k=|F|\ge2$ and $\mu_i=\mu_0$ for $i\in F$, where
$\mu_0>0$ by (H). All permutations of $F$ form a group $S$ as above, so
$c_F=c(\alpha^*_F)$, and $\lambda_F=d-k$ (Example~\ref{ex:models}). Both $f_T$
and $C_F$ are linear in the restriction of $\mu$ to $F$. So replacing $\mu$
by the uniform law on $[d]$ multiplies both by $1/(d\mu_0)$, and by
Theorem~\ref{thm:local} it does not change $c$. So we may take $\mu$ uniform.
Then the chain is the model of~\cite{paperB} with $p=1-(d-1)h$
(Example~\ref{ex:models}), and a vertex bin has probability $p^{N-1}/d$. Put
$s=k-1$, $z=Nh/p$ and $\mathcal D_k=k^{k+1/2}(4\pi)^{-s/2}$, and let $n$ be a balanced
bin with support $F$, so that $|n_i-N/k|<1$ for $i\in F$.
By~\cite[Eq.~(7)]{paperB}, which gives the ratio of $P_N(n)$ to $p^{N-1}/d$,
uniformly for $\frac12\log N\le z\le2\log N$,
\[
P_N(n)=\frac{p^{N-1}}d\,\mathcal D_k\,\frac{z^{s/2}e^{sz}}{N^s}
\Bigl(1-\frac s{8z}+O\Bigl(\frac1{z^2}+\frac{z^2}N\Bigr)\Bigr).
\]
Given a large $T$, let $N=\lceil e^T\rceil$. Then $T\le\log N<T+1$,
$T^2/N\le T^2e^{-T}$ and $z-T=(d-1)hz=O(T^2/N)$, so $z$ lies in the range
above. Also
$(N-1)\log p=-(d-1)T+O(T^2/N)$. Substituting, and using $d-1-s=d-k$, we get
\[
P_N(n)=N^{-s}\,\frac{\mathcal D_k}d\,T^{s/2}e^{-(d-k)T}\Bigl(1-\frac s{8T}
+O(T^{-2})\Bigr).
\]
On the other hand $|n'/N-\alpha^{*\prime}_F|\le k/N$, so $n/N\in U$ for large
$T$. We use Lemma~\ref{lem:discrete}, Theorem~\ref{thm:local} at $n/N$ and the
Taylor bounds, where $\nabla I_F(\alpha^*_F)=0$ gives
$T(I_F(n/N)-\lambda_F)=O(T/N^2)$. This gives
\[
P_N(n)=N^{-s}T^{s/2}K_Fe^{-(d-k)T}\Bigl(1+\frac{c(\alpha^*_F)}T
+O(T^{-3/2})\Bigr).
\]
Letting $T\to\infty$ in the ratio of the 2 expressions gives first
$K_F=\mathcal D_k/d$, in agreement with Corollary~\ref{cor:paperB}, and then
$c(\alpha^*_F)=-s/8$. So $c_F=-(k-1)/8$.
\end{proof}
